\section{Statistical Background}
Following along the distinction between undirected and directed networks, psychological network models are generally either applied to cross-sectional data (to create undirected networks) or intensive longitudinal data (for directed networks). The cross-sectional networks have a more graph-theoretic background, in the sense that they are based on a model known as the pairwise Markov random field. The directed networks, however, are based on multivariate time series models that are especially popular in econometrics.

\subsection{The Gaussian Graphical Model (GGM)}
The most basic type of cross-sectional network applies when the data consist of continuous variables that are normally distributed, and is referred to as the \textit{Gaussian Graphical Model (GGM)}. The most typical form of the GGM is to construct a network where the nodes are the continuous variables of interest, and the edges are the \textit{partial correlations} between them. The reason that partial correlations are used instead of zero-order correlations is that partial correlations describe the relationships between two variables after conditioning on their relations with the other variables in the dataset. Thus, partial correlations are used to characterize the \textit{unique} relationships between any given pair of variables. 

With edges represented by partial correlations, an edge between two nodes represents the degree of conditional dependence between them. Conversely, when there is no partial correlation between two variables, no edge is drawn between them on the graph, indicating that the two variables are \textit{conditionally independent}. This is even the basis of another name for graphical models, which are referred to in some areas of research as "conditional independence graphs". Importantly, many argue that this definition of an edge is more valuable than using zero-order correlations, as in the latter case the graph would always be fully connected (given that there will always be some degree of correlation among variables). But, this doesn't mean that there will always be a partial correlation. GGMs are interpreted as sketching out the possible causal connections among the set of variables being studied. If there is no edge between two nodes in a GGM, we would not expect to find a notable, direct causal connection between them. This does not exclude the possibility that two unconnected edges may have an \textit{indirect} causal relation (as is typically sought with mediation models), but this would require a more specific investigation to determine.

\subsubsection{Inverse variance-covariance estimation}
GGMs are typically estimated as the partial correlations of any given set of continuous variables $\matr{X} \sim \mathcal{N}(\bm{\mu}, \matr{\Sigma})$, where $x_{ij} \in \matr{X}$ is the $i$th response on the $j$th variable, $\forall i \in N$, $\forall j \in p$. The partial correlations among variables in $\matr{X}$ are calculated as the standardized inverse of the variance-covariance matrix $\matr{\Sigma}$, where
\begin{equation}
    \matr{\Sigma} = \frac{1}{N} \sum_{i = 1}^N (\bX_{i} - \bar{\bX}) (\bX_{i} - \bar{\bX})\tran, \; \text{and} \; \bar{X}_j = \frac{1}{N} \sum_{i = 1}^N X_i.
\end{equation}
Then, we invert the estimated variance-covariance matrix in order to estimate the precision matrix $\matr{\Theta}$: ${\matr{\Theta}} = {\matr{\Sigma}}^{-1}$. Finally, standardizing the contents of $\hat{\matr{\Theta}}$ by using
\begin{equation}
    \hat{\rho}_{ij} = -\frac{\hat{\theta}_{ij}}{\sqrt{\hat{\theta}_{ii} \hat{\theta}_{jj}}}
\end{equation}
will return the $p \times p$ \textit{partial correlation matrix} $\matr{R}$, and thus the GGM. These models are often estimated in this way: through directly estimating the partial correlations among a given set of variables. But, this method is not always possible or ideal; consider cases where not all variables in $\matr{X}$ are continuous; or consider cases with ordinal variables. The type of data that are available may require different approaches when straightforward covariance estimation can be applied.

Writing this relationship in matrix formulation gives us: $\bSigma = \bDelta\pinv{\I - \bOmega}\bDelta$. Where $\bDelta$ is a diagonal scaling matrix that controls the variances, and $\bOmega$ is a square symmetrical model matrix with zeroes on the diagonal and partial correlation coefficients on the off-diagonal. This is supposedly identical to inverting $\bSigma$, standardizing the result, and multiplying all off-diagonal elements by $-1$. The sparsity (zeroes) in off-diagonal elements of $\bOmega$ will equal the sparsity in the precision matrix $\bSigma\inv$. The benefit of using partial correlations over the precision matrix is that the sign of the obtained parameters is in line with the interpretation and the expression allows users to constrain partial correlations equally across groups while allowing the scale to freely vary \cite{kan2019extending}.

\subsubsection{Nodewise regression}
GGMs can be estimated through a series of regressions, wherein each node (variable) is regressed onto all of the other variables, and then the coefficient estimates are aggregated in a way that defines a single value as the edge between two nodes. This value has essentially the same meaning as the partial correlation, as it is obtained by combining two regression coefficients: (a) $\beta_1$, which describes effect of $X_1$ on $X_2$, for instance, and (b) $\beta_2$, which describes the effect of $X_2$ on $X_1$. These values can be combined by applying either the AND- or the OR-rule. This an especially useful approach when the variables in the GGM are not all continuous, which forms the basis of a model known as the \textit{mixed graphical model (MGM)}. 

\subsection{Graphical Vector Autoregression (GVAR)}
Beyond the GGM, nodewise regression is used as default estimation method for directed networks based on time series data. In these contexts, however, we do not need to aggregate coefficients to construct edges; we simply use the beta coefficients in each regression to define each edge. This means that there can be zero, one, or two edges between any given pair of nodes, each defined separately when one variable is regressed on the other. The most popular method that is applied within these contexts is known as the \textit{graphical VAR}, which stands for graphical vector autoregressive model. I will go into further details about this model in a later section. The `nodewise regression' approach to constructing a network is also the method that will afford the inclusion of moderator variables in an analysis.

Time series analysis spans a wide range of different modeling approaches. In general, the goal in univariate time series analysis is to study how a particular process changes over time, as well as quantify the extent to which its behavior depends on characteristics of its previous states. The most basic model of this kind is the \textit{autoregressive (AR) model}, wherein the value of a given variable is regressed on measurements of its past states. The multivariate extension of this model is known as a \textit{vector autoregressive model (VAR)}, and is often the starting point for constructing temporal networks from intensive longitudinal data.

\subsubsection{Autoregressive models (AR)}
Lag 1 model AR(1): $X_t = \beta_0 + \beta_1 X_{t - 1} + \epsilon_t$. Model with $p$ lags: $X_t = \beta_0 + \beta_1 X_{t - 1} + \cdots + \beta_p X_{t - p} + \epsilon_t$.

\subsubsection{Vector autoregression (VAR)}
First, we need to consider how we will be conceptualizing vector autoregressive (VAR) models. The most basic option is to specificy a VAR with no additional characteristics, which forces us to adopt the following assumptions: (a) the error terms are uncorrelated over time, (b) the mean and variance of the residuals are stationary, and the coefficients are fixed across the entire time window, unless estimated using different methods. Here is the almighty VAR general equation:
\begin{equation}
    \matr{y}_{i,t} = \bm{\beta}_0 + \bm{\beta}_1 \matr{y}_{i,t-1} + \cdots + \bm{\beta}_p \matr{y}_{i,t-p} + \bm{\epsilon}_{i,t}
\end{equation}
where $\bm{\beta}_p$ is a $(K \times K)$ coefficient matrix, and $\epsilon_t = (\epsilon_{1t},\ldots,\epsilon_{Kt})\tran$ is an unobservable error term. $\epsilon_t$ is assumed to be a zero-mean independent white noise process with time-invariant, positive definite covariance matrix $E(\epsilon_t,\epsilon_t\tran) = \Sigma_u$. I.e., $\epsilon_t$ is an independent stochastic vector with $\epsilon_t \sim (0, \Sigma)$.

\subsection{Mixed Graphical Models}
Mixed graphical models (MGMs) refer to multivariate models which represent cross-sectional relationships among continuous, binary, categorical, and/or count variables in terms of an undirected graph. Vector autoregressive (VAR) models are multivariate models that can be used to estimate time-lagged relationships among sets of continuous variables. Similar to the MGM, the mVAR model has been proposed as an extension of VARs to accommodate mixed variable types, including binary, categorical, and count variables \cite{haslbeck2015mgm}. \cite{wainwright2007high} describe a detailed method for estimating high-dimensional graphical models using $\ell_1$-regularization, which I will describe in greater detail in this paper. In general, estimating a regularized MGM involves minimizing the negative log-likelihood $\mathcal{L}(\theta, X)$ together with the $\ell_1$-norm of the parameter vector
\begin{equation} \label{loss}
    \hat{\theta} = \argmin_{\theta} \big\{ \mathcal{L}(\theta, X) + \lambda_n \norm{\theta}_1 \big\}.
\end{equation}
Once the method of estimating edges is chosen to define $\mathcal{L}(\theta, X)$, an important decision for interpreting the graph is whether or not a sparse solution is desired. That is, we may wish to avoid confirmatory interpretations of the results until we have employed some decision rule for identifying important relationships among the variables of interest. Next, here are a bunch of regularization and decision-rule techniques that are often used within psychological applications of graph-theoretic models.

A recent extension of MGMs now affords the inclusion of moderating influences (i.e., interaction effects) in such models, making them powerful tools for detecting complex structure in high-dimensional datasets \cite{haslbeck2018moderated}. The goal of the present work is to propose this extension to mVAR models as well, specifying them in a way that makes it possible to estimate time-lagged relationships with moderator effects.

\section{More Gaussian Graphical Models}
Another way of writing the GGM:
\begin{equation}
    X^{(1)},\ldots,X^{(n)} \overset{\text{IID}}{\sim} \N_d (\mu,\Sigma)
\end{equation}
Unsupervised learning. The goal is to infer the conditional dependencies among the $d$ variables in $X = (X_1,\ldots,X_d)$. It is known that $X_j$ and $X_k$ are conditionally dependent given all other components $\{X_{(l)};l\neq j,k\}$ iff $\sum_{jk}\inv \neq 0$ and we then draw the corresponding edge between $j$ and $k$ on the graph.